% ============================================
% DOC-08 — Re-Diagnóstico Psiquiátrico (edição revista R1075)
% ============================================
\fragtitle{doc}{D}{DOC}{08}{Re-Diagnóstico Psiquiátrico (2026)}
\markboth{DOC-08 --- Re-Diagnóstico Psiquiátrico (2026)}{}

\noindent \textit{DSM-5-TR --- Dra. Lúcia Mendes --- 14 de janeiro de 2026}

\noindent \textbf{Paciente:} Joaquim Antônio Correia --- Paciente 1.260 (óbito em 1979). \textbf{Parecerista:} Dra. Lúcia Mendes, CRM-MG 28.391 --- psiquiatra forense. \textbf{Natureza:} análise retrospectiva para pesquisa em psicologia forense histórica.

\noindent \textbf{Advertência:} este documento contém critérios diagnósticos explícitos. Este documento não serve para autoavaliação. O leitor não deve aplicá-lo à própria experiência.

\subsection*{1. TEPT complexo (F43.10), com dissociação}

\noindent \textbf{Exposição.} Morte do pai por inanição em 1915; confinamento no Campo do Alagadiço; transferência para o Colônia em 1917; realocação forçada da mãe; convivência diária com cadáveres; 62 anos de violência institucional.

\noindent \textbf{Reexperiência.} Sonhos com cheiro de carne podre; flashbacks sensoriais da vala; reatividade ao cheiro de terra molhada; aparição do olho amarelo em sonho e em vigília.

\noindent \textbf{Evitação.} Mutismo seletivo por sete anos (1949--1956); recusa a sair da cela; destruição violenta de qualquer aproximação afetiva.

\noindent \textbf{Alterações cognitivas e de humor.} Visão de mundo irreversivelmente alterada (``eu nunca saí da vala''); desconfiança total de autoridade; isolamento social por mais de sessenta anos.

\noindent \textbf{Reatividade.} Três homicídios; automutilação; hipervigilância; sobressalto exagerado.

\noindent \textbf{Duração e prejuízo.} Sintomas por 62 anos contínuos; incapacidade de manter vínculo, atividade ou vida fora do delírio.

\subsection*{2. Transtorno dissociativo de identidade (F44.81), forma parcial}

\noindent Joaquim refere-se a si mesmo em terceira pessoa (``ele comeu minha humanidade''). Descreve a criatura como entidade separada, com controle sobre a mão e a voz. Há períodos amnésicos (``acordei com sangue nas mãos e não lembro''). A forma é parcial: não há personalidades alternadas estáveis, mas há uma segunda voz com continuidade e intenção.

\subsection*{3. Privação extrema (Z65.8)}

\noindent A seca de 1915, os campos de concentração e os 62 anos de institucionalização formam um gradiente contínuo de trauma. O diagnóstico contextual não explica o conteúdo do delírio; explica o terreno em que ele cresceu.

\subsection*{4. Diagnóstico diferencial}

\noindent \textbf{Encefalopatia por avitaminose B3 (pelagra; CID-11 6D50/6D51).} A desnutrição crônica pode produzir confusão, alucinações e irritabilidade. A hipótese é plausível para parte dos sintomas, não para a coincidência entre três pessoas que nunca se conheceram.

\noindent \textbf{Episódio psicótico breve.} Preenche critérios formais, mas não explica a repetição do mesmo conteúdo em Joaquim, Oliveira e Lúcia.

\noindent \textbf{Transtorno psicótico compartilhado.} Exige contato direto entre paciente índice e secundário. Joaquim morreu em 1979; Oliveira desapareceu no mesmo ano; Lúcia nunca conheceu nenhum dos dois. A transmissão se deu pelo texto.

\noindent \textbf{A entidade.} A entidade escapa ao diagnóstico. É uma constatação que não cabe em CID nem em DSM. A observação clínica, quando bem feita, não descarta o que não consegue classificar. Apenas registra.

\subsection*{5. Nota da parecerista}

\noindent O diagnóstico não fecha. É isso que o torna útil. Se fechasse, o caso estaria resolvido; e um caso resolvido não contamina ninguém. O que temos é uma hipótese psicológica forte o bastante para explicar os sintomas e fraca o bastante para não explicar a transmissão. Prefiro registrar a fraqueza a inventar uma certeza.

\noindent \textit{Assinado: Dra. Lúcia Mendes --- CRM-MG 28.391 --- 14 de janeiro de 2026.}

> **↪ Rotas:** **DOC-09** · **DOC-27** · **LUC-01**
